\setcounter{theorem}{16}
	\subsection{Pseudorandomness}
	
	A PRG against a class of functions $\mathfrak{C}$ is a deterministic procedure $G$ mapping short Boolean strings (seeds) to longer Boolean strings, so that $G$'s output ``looks random" to every function in $\mathfrak{C}$.
	
    \begin{definition}[Pseudorandom generators]
    	Let $G\colon\boolpm^{\ell}\to\boolpm^n$ be a function, $\mathfrak{C}$ be a class of Boolean functions, and $0<\varepsilon<1$. We say that $G$ is a \emph{pseudorandom generator of seed length $\ell$ that $\varepsilon$-fools $\mathfrak{C}$} if, for every function $f\in\mathfrak{C}$, it is the case that
    	\[
    	    \abs{\Exp_{z\sim\boolpm^\ell}\csqbra{f\cprn{G\cprn{z}}}-
    	    \Exp_{x\sim\boolpm^n}\csqbra{f\cprn{x}}}\leq\varepsilon.
    	\]
    \end{definition}
	
	A PRG $G$ outputting $n$ bits is called \emph{explicit} if $G$ can be computed in $\poly(n)$ time. All PRGs stated in this paper are explicit.
	
